\section{Introduction}\label{sec:intro}

A circuit calculus is a programming language of a particular kind: its
programs are circuits, words of gates on numbered wires, and its
equational theory is a finite list of circuit equations.  Such a list is
\emph{sound} when each equation holds in the intended semantics and
\emph{complete} when every semantic equality between circuits is
derivable from it.  Completeness is what a language designer wants from
an equational theory: it licenses normal-form-based decision procedures
for program equivalence, and it tells the author of a circuit optimizer
that no valid rewrite is missing.  It is also, for quantum circuits, the
kind of theorem that resists being trusted on paper.  Soundness is a
finite collection of matrix identities.  Completeness is proved by a
\emph{normal form}: one shows that every circuit rewrites to a normal
form using the equations, and that distinct normal forms denote distinct
operators.  The middle step is a case analysis, every generator against
every piece of every normal form, that is mechanical, unforgiving and
enormous.  Selinger's presentation of the $n$-qubit Clifford
group~\cite{selinger2015clifford} runs a substantial rewrite system
through such an analysis; its qutrit successor~\cite{li2025qutrit} needs
more than fifty pages; and its generalisation to all odd primes by Bian,
Li, Ross, van de Wetering, and Zhao~\cite{bian2026qudit} --- sixteen
rewrite rules for multi-qudit Clifford circuits in every odd prime
dimension $p$, henceforth ``the qupit paper'' --- rests on roughly a
hundred pages of derivations spread over its appendices and a
supplement~\cite{bian2026supplement}.

This paper reports a machine-checked proof of that theorem, and, more
importantly for a programming-languages audience, a \emph{method} for
machine-checking theorems of this shape.  We formalise, in safe,
axiom-free Agda~\cite{norell2007thesis,agdastdlib}, that the qupit
paper's rule set presents the multi-qupit Clifford group at every width
$n$ and every odd prime $p$, against a semantics constructed
independently of the syntax.  The formalisation is built on a library
for presentations, normal forms and completeness of circuit relation
sets~\cite{bian2026framework,qupit-code}, whose design --- inductively
defined circuits, presented monoids as inductive congruences, normal
forms as function bundles, a verified Reidemeister--Schreier engine, and
presentation theorems for products and extensions --- we recall in
\S\ref{sec:framework}.  Every normal form in the development is an
executable function, so each completeness theorem doubles as a verified
decision procedure for circuit equivalence in its calculus.

\paragraph{The method.}
The formalisation does \emph{not} follow the qupit paper's proof, and
the reason is the paper's main methodological lesson.  The paper works
with a single rule set for the whole Clifford group and organises its
completeness proof around the symplectic group $\Sp$, the Clifford
group modulo Paulis: it shows that every circuit rewrites, \emph{up to
Paulis}, to a symplectic normal form, and then boosts symplectic
completeness to projective and unitary completeness by a general
argument (its Proposition~4.9) in which each rewrite is repaired by
pushing a Pauli correction to the end of the circuit.  Throughout,
relations are used ``projectively'': Paulis and scalars are silently
dropped and later restored.  This is sound mathematics, and we believe
it; but it is a poor specification for a verified normalisation
procedure.  The rewriting steps preserve no single syntactic invariant,
the correction Paulis are not part of the data the rewrites compute
with, and the termination and well-definedness arguments are stated at a
level --- ``push the Pauli to the right, it changes to some other
Pauli'' --- that a proof assistant would have to reconstruct from
scratch.

We therefore chose a route that a proof assistant likes: \emph{factor
the group, present the factors, compose the presentations}.  The
projective Clifford group is a semidirect product of the projective
Pauli group by the symplectic group, and the Clifford group is a central
extension of that by the scalars.  We present the three factors
separately --- the symplectic group by a coset-table normal form, the
Pauli group compositionally, the scalars as a cyclic group --- and glue
the presentations together with generic, verified presentation theorems
for semidirect products and central extensions.  On this route the
Paulis are never implicit: the symplectic development never mentions
them, and they enter only through the conjugation action of the
semidirect product, which is verified once, by semantics, rather than
threaded through every rewrite.  A final layer of presentation-level
``plumbing'' reduces the glued alphabet and relations to the paper's
generators $H$, $S$, $CZ$ and its rules, and a verified isomorphism of
presentations connects the two routes: the paper's Figure~1 read modulo
scalars presents the same group as our semidirect presentation.

\subsection{Contributions}\label{sec:contributions}

\begin{enumerate}
\item[(1)] \textbf{A method, and a comparison.}  We explain precisely
  where the qupit paper's proof resists formalisation and how the
  factor-present-compose route avoids each difficulty
  (\S\ref{sec:two-routes}), introduce the route on two warm-up examples,
  the symmetric groups and the wreath products $\mathbb{Z}_N \wr S_n$
  (\S\ref{sec:warmup}), and report what the formalisation covers and
  what it does not (\S\ref{sec:discussion}).  The route is generic: its
  two presentation theorems were proved for arbitrary groups before this
  project began.
\item[(2)] \textbf{A machine-checked presentation of the projective
  qupit Clifford groups.}  For every odd prime $p$ and every width $n$,
  the paper's Figure~1 read modulo scalars presents
  $\PPauli \rtimes \Sp$ (Theorem~\ref{thm:main}), via a chain of
  verified isomorphisms of presentations from a semidirect presentation
  on $X, Z, H, S, CZ$ down to the paper's fifteen rules on $H, S, CZ$
  (\S\ref{sec:composing}).
\item[(3)] \textbf{A machine-checked presentation of $\Sp$ and a
  verified unique normal form.}  Eighteen relations present the
  symplectic groups; the normal form is the paper's, refined into a
  doubly inductive datatype of \emph{normal boxes} and \emph{rows},
  realised as a Reidemeister--Schreier coset tower, and proved unique
  against the symplectic action on Pauli vectors
  (Theorem~\ref{thm:unique}, \S\ref{sec:symplectic}).
\item[(4)] \textbf{The relation reductions, verified.}  The box
  relations that the normalisation uses are derived from the full
  symplectic rules, those from the simplified ones, the semidirect
  relations from the Clifford ones, and the paper's rules from ours;
  each reduction is a checked isomorphism of presentations
  (\S\ref{sec:reduction}, \S\ref{sec:plumbing}).
\item[(5)] \textbf{The scalars restored.}  The exact rule set --- the
  sixteen rules corrected by the phase each actually picks up, a single
  power of $\omega$ on one rule --- presents a central extension of the
  projective group by $\langle\omega\rangle$, verified against the Weyl
  cocycle (\S\ref{sec:scalars}).
\end{enumerate}

Everything stated as a theorem in this paper typechecks under
\texttt{--safe}: there are no postulates, no holes and no classical
axioms.  Two facts are \emph{not} formalised, and we say so wherever
they are used: that the projective Clifford group of unitaries is
$\PPauli \rtimes \Sp$, and that the Clifford group is a central
extension of it by $\langle\omega\rangle$.  Both are
classical~\cite{bian2026qudit}, neither involves circuits, and
\S\ref{sec:discussion} states exactly what our semantic groups are.

\paragraph{Outline.}
Section~\ref{sec:background} recalls qupit Clifford operators and states
the two main theorems as Agda types.  Section~\ref{sec:two-routes}
contrasts the paper's proof with ours.  Section~\ref{sec:framework}
recalls the framework and \S\ref{sec:warmup} exercises it on the two
warm-up examples.  Sections~\ref{sec:symplectic}
and~\ref{sec:composing} are the method: the symplectic presentation,
then the composition of the three factors.  Section~\ref{sec:discussion}
gives statistics, engineering lessons and caveats; \S\ref{sec:related}
surveys related work; \S\ref{sec:conclusion} concludes.
